\begin{ex}
Số gồm 5 chục nghìn, 3 nghìn, 6 trăm, 2 chục và 8 đơn vị viết là gì?
\shortans{53628}
\loigiai{
Số đó là 53628.
}
\end{ex}

\begin{ex}
Số 74152 có chữ số hàng chục nghìn là bao nhiêu?
\shortans{7}
\loigiai{
Chữ số hàng chục nghìn là 7.
}
\end{ex}

\begin{ex}
Số 80425 có chữ số hàng nghìn là bao nhiêu?
\shortans{0}
\loigiai{
Chữ số hàng nghìn là 0.
}
\end{ex}

\begin{ex}
Số liền sau của số 99 999 là số nào?
\shortans{100000}
\loigiai{
Số liền sau của 99 999 là 100 000.
}
\end{ex}

\begin{ex}
Số liền trước của số 50 000 là số nào?
\shortans{49999}
\loigiai{
$50000 - 1 = 49999$.
}
\end{ex}

\begin{ex}
Số lớn nhất có năm chữ số là số nào?
\shortans{99999}
\loigiai{
Số lớn nhất có 5 chữ số là 99999.
}
\end{ex}

\begin{ex}
Số bé nhất có năm chữ số là số nào?
\shortans{10000}
\loigiai{
Số bé nhất có 5 chữ số là 10000.
}
\end{ex}

\begin{ex}
Số tròn chục nghìn liền sau của 60 000 là số nào?
\shortans{70000}
\loigiai{
Số tròn chục nghìn liền sau là 70 000.
}
\end{ex}

\begin{ex}
Số tròn nghìn liền trước của 45 000 là số nào?
\shortans{44000}
\loigiai{
Số tròn nghìn liền trước là 44 000.
}
\end{ex}

\begin{ex}
Số tròn chục liền sau của 81 240 là số nào?
\shortans{81250}
\loigiai{
Số đó là 81250.
}
\end{ex}

\begin{ex}
Viết số thích hợp: $60000 + 4000 + 300 + 20 + 5 = \dots$
\shortans{64325}
\loigiai{
Số đó là 64325.
}
\end{ex}

\begin{ex}
Viết số thích hợp: $80000 + 500 + 40 + 7 = \dots$
\shortans{80547}
\loigiai{
Số đó là 80547.
}
\end{ex}

\begin{ex}
Viết số thích hợp: $90000 + 8000 + 6 = \dots$
\shortans{98006}
\loigiai{
Số đó là 98006.
}
\end{ex}

\begin{ex}
Số 75 412 có giá trị của chữ số 7 là bao nhiêu?
\shortans{70000}
\loigiai{
Chữ số 7 ở hàng chục nghìn có giá trị 70 000.
}
\end{ex}

\begin{ex}
Số 38 650 có giá trị của chữ số 8 là bao nhiêu?
\shortans{8000}
\loigiai{
Chữ số 8 ở hàng nghìn có giá trị 8000.
}
\end{ex}

\begin{ex}
Cho các chữ số 4, 1, 9, 0, 7. Số lớn nhất có năm chữ số khác nhau lập từ các chữ số trên là số nào?
\shortans{97410}
\loigiai{
Sắp xếp giảm dần: 97410.
}
\end{ex}

\begin{ex}
Cho các chữ số 5, 2, 0, 8, 3. Số bé nhất có năm chữ số khác nhau lập từ các chữ số trên là số nào?
\shortans{20358}
\loigiai{
Chữ số đầu khác 0 nhỏ nhất là 2, tiếp theo là 0, 3, 5, 8: 20358.
}
\end{ex}

\begin{ex}
Có bao nhiêu số tròn chục nghìn từ 10 000 đến 100 000?
\shortans{10}
\loigiai{
10000, 20000, ..., 100000 (10 số).
}
\end{ex}

\begin{ex}
Điền số thích hợp: $69999 < \dots < 70001$.
\shortans{70000}
\loigiai{
Số ở giữa là 70000.
}
\end{ex}

\begin{ex}
Số liền sau của số 45 899 là số nào?
\shortans{45900}
\loigiai{
$45899 + 1 = 45900$.
}
\end{ex}

\begin{ex}
Số liền trước của số 32 000 là số nào?
\shortans{31999}
\loigiai{
$32000 - 1 = 31999$.
}
\end{ex}

\begin{ex}
Hiệu của số bé nhất có năm chữ số và số lớn nhất có bốn chữ số là bao nhiêu?
\shortans{1}
\loigiai{
$10000 - 9999 = 1$.
}
\end{ex}

\begin{ex}
Số 88 888 có bao nhiêu chữ số 8?
\shortans{5}
\loigiai{
Có 5 chữ số 8.
}
\end{ex}

\begin{ex}
Trong số 65 321, chữ số hàng trăm là chữ số nào?
\shortans{3}
\loigiai{
Chữ số hàng trăm là 3.
}
\end{ex}

\begin{ex}
Giá trị của chữ số 5 ở hàng chục nghìn gấp mấy lần giá trị của chữ số 5 ở hàng trăm?
\shortans{100}
\loigiai{
$50000 : 500 = 100$ lần.
}
\end{ex}
